\documentclass{article}
\usepackage{amsmath,amssymb}
\begin{document}
\title{Nonvanishing of Dirichlet \(L\)-functions Right of One Half}
\author{}
\date{}
\maketitle

\section*{Main line}
For every positive modulus \(q\), character \(\chi\bmod q\), and \(\rho\) with \(\Re\rho>1/2\), the target asserts \(L(\rho,\chi)\ne0\), with the principal character at \(\rho=1\) excepted. Put
\[
 M=18480q,\quad a=9240q-1,\quad
 \varepsilon=\frac{\min(\Re\rho,1)-1/2}{51},\quad
 \beta=\tfrac12+8\varepsilon.
\]
Then \(0<\varepsilon<1/100\) and \(1/2<\beta<\Re\rho\). Fix a prime \(d\equiv a\pmod M\). It satisfies the hypotheses of Gap 1. For large \(x\), take its witnesses and apply Gap 2 with \(t=10^6\log x\) and \(c=T/H\). Its \(u\) can be used in Gap 1, giving
\[
 R_{d,M}(x)\le
 \frac{f(\log b)^5}{mH}+x^{-10^6}
 \le \frac{\varphi(M)x^{6\varepsilon}}{\sqrt{x}}+x^{-10^6}
 \le x^{\beta-1}
\]
eventually. The allowed one-sided criterion makes the entire \(L\)-function of the character induced to \(M\) nonzero at \(\rho\). The allowed finite Euler-factor fact transfers this to \(L(\rho,\chi)\), subject to the stated exception.

\section*{Gaps to prove independently}
For a prime \(d\equiv7\pmod8\) coprime to \(N\), let \(\kappa_d(n)=(n/d)\), let \(\Lambda\) be the von Mangoldt function, and set
\[
\begin{split}
 R_{d,N}(y)
 &=\sum_{1\le n\le y}\Lambda(n)(1-\kappa_d(n))
   \mathbf1_{n\equiv1\;(\mathrm{mod}\;N)}
   \left(\frac1n-\frac1y\right)
   -\frac{\log y}{\varphi(N)}-C_{d,N},\\
 C_{d,N}
 &=\Re\!\left[
 \frac{\displaystyle\sum_{\eta\bmod N}
 \left(\frac{\widetilde L'_{\eta\kappa_d}(1)}
 {\widetilde L_{\eta\kappa_d}(1)}
 -\frac{\widetilde L'_\eta(1)}{\widetilde L_\eta(1)}\right)-1}
 {\varphi(N)}\right].
\end{split}
\]
Here \(\widetilde L_\eta(s)=(s-1)L(s,\eta)\) for a principal character and \(L(s,\eta)\) otherwise.

\noindent\textbf{Gap 1.} Let \(M>0\) be divisible by \(18480\), let \(d\) be prime with \(M\mid2(d+1)\) and \(M\nmid d+1\), and let \(0<\varepsilon<1/100\). For every sufficiently large \(x\), there are positive integers \(m,H,T,f\) and positive reals \(a,b\) such that
\[
 m\ge\sqrt{x},\quad \log b=x^\varepsilon,\quad
 H\ge T/\log b,\quad f=T\varphi(M),
\]
and, for \emph{every} \(u\in[10^6\log x,(10^6\log x)^2]\),
\[
 maH\left(R_{d,M}(x)-R_{d,M}(e^u)
              +\frac TH R_{d,1}(e^u)\right)
 \le fa(\log b)^5.
\]
This is the geometric step used above. Construct the cyclotomic tower and its thick places, then compare upper and product-formula lower bounds for the same covolume of theta lifts. Their common main height cancels; the remaining prime sum gives the displayed inequality. The construction must allow the auxiliary scale \(e^u\) throughout the stated interval.

\medskip
\noindent\textbf{Gap 2.} Let \(d\equiv7\pmod8\) be prime and coprime to \(N>0\). For every sufficiently large \(t\) and every \(0\le c\le e^t\), there is \(u\in[t,t^2]\) such that
\[
 R_{d,N}(e^u)-cR_{d,1}(e^u)\le e^{-t}.
\]
This supplies the auxiliary scale in the main line. Prove it from the Mellin transforms of the two errors by a localized moment argument.

\section*{Allowed facts}
The one-sided prime-error nonvanishing criterion, finite Euler-factor behavior when changing modulus, and the elementary arithmetic of \(9240q-1\pmod{18480q}\).

\end{document}
